\chapter{The generative law: first-passage finishing times}
\label{ch:gen}
\label{ch:generative}

\section{From a velocity profile to a finishing time}

Fix a race of geometric length $L$ (metres). A horse covers the trip by accumulating distance at
a time-varying velocity $v_i(t)$, so its finishing time is the first time its accumulated
distance reaches $L$:
\begin{equation}
  \Tfin_i \;=\; \inf\Big\{\, t \ge 0 : \textstyle\int_0^t v_i(s)\,ds \;\ge\; L \,\Big\}.
  \label{eq:firstpassage-def}
\end{equation}
This is already a first-passage statement: the finishing time is the hitting time of a level $L$
by the integrated-velocity process. Everything in the theory of racing is, in the end, a theory
of \emph{what sets $v_i(t)$} --- how much speed the engine can produce, how that capacity is
deployed across the trip, how the arena and the rival field perturb it, and what is irreducibly
random. The chapters that follow populate $v_i(t)$ with mechanism. The present chapter fixes the
\emph{stochastic} structure of \eqref{eq:firstpassage-def} that turns a per-horse velocity model
into win probabilities.

\section{The drift--diffusion idealisation and the Wald law}

The production model idealises the accumulated-distance process $X_i(t) = \int_0^t v_i\,ds$ as a
Brownian motion with drift: in increments,
\begin{equation}
  dX_i \;=\; v_i\, dt \;+\; \sqrt{2\Ddiff_i}\; dW_i,
  \qquad X_i(0)=0,
  \label{eq:sde}
\end{equation}
where $v_i>0$ is the (effective, trip-averaged) drift velocity, $\Ddiff_i>0$ is a per-horse
diffusion coefficient with units of $\mathrm{m^2/s}$, and $W_i$ is a standard Wiener process. The
finishing time is the first passage of $X_i$ to the absorbing level $L$,
$\Tfin_i = \inf\{t : X_i(t)=L\}$. It is a classical result that the first-passage time of a
drift--diffusion to a fixed positive level follows the \emph{Inverse-Gaussian} (Wald)
distribution,
\begin{equation}
  \Tfin_i \;\sim\; \IG(\mu_i,\lambda_i),
  \qquad
  \mu_i \;=\; \frac{L}{v_i},
  \qquad
  \lambda_i \;=\; \frac{L^2}{2\Ddiff_i},
  \label{eq:wald-params}
\end{equation}
with density, for $t>0$,
\begin{equation}
  f_{\IG}(t;\mu,\lambda)
   \;=\; \sqrt{\frac{\lambda}{2\pi t^{3}}}\;
   \exp\!\left(-\,\frac{\lambda\,(t-\mu)^2}{2\mu^{2}\,t}\right).
  \label{eq:ig-density}
\end{equation}

Two derivation checks make the parameterisation transparent. First, the mean: the expected
finishing time is the distance divided by the drift,
$\E[\Tfin_i] = \mu_i = L/v_i$, exactly the deterministic ``run the trip at speed $v_i$'' answer.
Second, the variance. For the Inverse-Gaussian, $\Var(\Tfin_i)=\mu_i^3/\lambda_i$, so
substituting \eqref{eq:wald-params},
\begin{equation}
  \Var(\Tfin_i)
   \;=\; \frac{\mu_i^3}{\lambda_i}
   \;=\; \frac{(L/v_i)^3}{L^2/(2\Ddiff_i)}
   \;=\; \frac{2\Ddiff_i\,L}{v_i^{3}}.
  \label{eq:wald-var}
\end{equation}
The variance of the finishing time is set by the diffusion $\Ddiff_i$ and \emph{grows linearly in
the trip length} $L$. This single fact has large downstream consequences (Section~\ref{sec:cvlaw}
and Chapter~\ref{ch:theory}): it predicts, with \emph{no free parameter}, that longer races are
intrinsically more predictable.

\begin{remark}[Why first-passage, not additive Gaussian]
One could model $\Tfin_i$ directly as $\mu_i + \varepsilon_i$ with Gaussian $\varepsilon_i$. That
is rejected on physical grounds: finishing times are strictly positive and right-skewed (a horse
can lose far more time to trouble than it can ever gain), and the first-passage construction
delivers exactly that --- the Inverse-Gaussian is supported on $(0,\infty)$ and right-skewed,
with the skew arising mechanically from the absorbing boundary. Empirically the additive-Gaussian
``\code{time}'' surrogate scores worse than the Wald law, and the Wald law ties or beats the
log-normal (velocity-multiplicative) surrogate while resting on a single physical diffusion
constant.
\end{remark}

\section{The three knobs}
\label{sec:threeknobs}

Equations \eqref{eq:wald-params}--\eqref{eq:wald-var} expose precisely three places a covariate
can act, and the entire feature programme is organised around them.

\begin{enumerate}[label=\textbf{Knob \arabic*.},leftmargin=*]
  \item \textbf{The drift $v_i$ (aerobic ceiling $\vc$).} Through $\mu_i=L/v_i$, anything that
        raises the sustainable speed lowers the expected time. This is the channel of raw ability:
        form, rating, sectional critical velocity, opponent-adjusted collateral form, speed
        figures, barrier trials, mass. The drift is built up in
        Chapters~\ref{ch:pop}--\ref{ch:competition} and assembled in Chapter~\ref{ch:assembly}.
  \item \textbf{The reserve $\Wp$ (a late-race boost to the effective drift).} The anaerobic
        reserve shortens the effective trip the aerobic engine must cover: a horse with reserve
        $\Wp$ behaves as if running a distance $L$ but with an effective drift scaled up by
        $(1+\Wp/L)$, so that
        \begin{equation}
          \mu_i \;=\; \frac{\Leff}{v_i\,\bigl(1+\Wp_i/L\bigr)},
          \label{eq:mu-wprime}
        \end{equation}
        where $\Leff$ is the geometric length corrected for ground-loss (Chapter~\ref{ch:geometry}).
        The factor $(1+\Wp/L)$ is the production form of Hill's critical-velocity law
        $T=(L-\Wp)/v_c$ (Chapter~\ref{ch:energetics}): the reserve buys back distance, and
        because it is divided by $L$ it \emph{bites harder over sprints than over staying trips},
        which is exactly the distance-aware behaviour the physiology demands. Reserve phenomena
        --- late-race fade resistance, the positional-fade cost a front-runner pays for setting
        the pace --- load this knob (Chapter~\ref{ch:assembly}, the ``T1'' recoupling).
  \item \textbf{The diffusion $\Ddiff_i$ (consistency / racing luck).} Through
        $\lambda_i=L^2/2\Ddiff_i$, the diffusion sets the spread of a horse's time around its
        mean. Because the winner is an \emph{order statistic}, variance is as load-bearing as the
        mean: a high-$\Ddiff$ horse has a fat left tail and upsets more often than its drift
        implies; a metronomic horse upsets less. This is the one knob the model exploited last and
        least; the draw$\to\Ddiff$ trip-variance term (Chapter~\ref{ch:geometry}) and the
        surface-dependent diffusion are its current uses.
\end{enumerate}

\begin{principle}[Right signal, right knob]
A feature is correctly structured only when it loads the knob its mechanism acts on. Loading a
reserve phenomenon (late fade) onto the aerobic ceiling $\vc$ is a \emph{correlated
approximation} --- it raises or lowers the whole velocity curve instead of just the closing kick
--- and is strictly dominated, where the data can resolve it, by loading $\Wp$. Likewise a
consistency signal belongs on $\Ddiff$, not on the mean.
\end{principle}

\section{The closed-form Plackett--Luce surrogate}
\label{sec:pl}

For very wide parameter sweeps it is convenient to have a sampling-free win-probability map. If
the finishing-time noise is taken to be Gumbel rather than Inverse-Gaussian, the argmin in
\eqref{eq:argmin} has the closed-form multinomial-logit (Plackett--Luce) solution. Writing the
deterministic mean times $\mu_i$ and a scale $\sigma_t$, the win probability is the softmax of the
negative mean times,
\begin{equation}
  p_i \;=\; \frac{\exp(-\mu_i/\sigma_t)}{\sum_{j}\exp(-\mu_j/\sigma_t)} ,
  \qquad \mu_i = \frac{\Leff}{v_i(1+\Wp_i/L)} .
  \label{eq:pl}
\end{equation}
Plackett--Luce additionally hands back the \emph{entire} finishing-order distribution in closed
form (the probability of any permutation factorises as a sequence of softmaxes over the remaining
field), which makes it the natural probe for place and exotic structure. It is, however, the
\emph{free / non-interacting} theory: independent finishing times, no traffic, no draft, no
tactical coupling. The production model uses the Wald Monte-Carlo (Chapter~\ref{ch:assembly}); PL
is retained as the instant scout and as the conceptual baseline against which the interaction
programme of Chapter~\ref{ch:theory} is defined.

\begin{remark}[The surrogate is faithful for win, by construction]
Interactions (draw, mass, traffic) have repeatedly tested as low-leverage for the \emph{win}
order statistic specifically --- the knife-edge for first place is dominated by the front of the
field, where the mean-field treatment is adequate. The pure-\code{numpy} Wald/PL surrogates
therefore match the full interacting Monte-Carlo on win to within noise while running roughly
$35\times$ faster, which is what makes the diagnose/recalibrate loop of Chapter~\ref{ch:method}
tractable. The interacting simulator is reserved for the place/exotic regime, where whole-field
ordering --- and hence the many-body terms --- matters.
\end{remark}

\section{The distance law of variance}
\label{sec:cvlaw}

Combine the mean $\mu_i=L/v_i$ and the variance \eqref{eq:wald-var} into the coefficient of
variation of the finishing time:
\begin{equation}
  \mathrm{CV}(\Tfin_i)
   \;=\; \frac{\sqrt{\Var(\Tfin_i)}}{\E[\Tfin_i]}
   \;=\; \frac{\sqrt{2\Ddiff_i L/v_i^3}}{L/v_i}
   \;=\; \sqrt{\frac{2\Ddiff_i}{L\,v_i}}
   \;\propto\; \frac{1}{\sqrt{L}} .
  \label{eq:cvlaw}
\end{equation}
The relative spread of finishing times \emph{falls as one over the square root of the trip
length}. Favourites should win more, and upsets should be rarer, over longer races --- a
prediction that contains no fitted parameter, falling out purely from the geometry of
first-passage. Chapter~\ref{ch:theory} develops the refinement (a distance-independent variance
floor injected by the start and by traffic, which would make sprints \emph{excess}-dispersed
relative to \eqref{eq:cvlaw}) and the order-statistic corollary that the $\Ddiff$ knob has more
leverage in larger fields. This prediction is, at the time of writing, under empirical validation
(queue item \#33).

\section{Summary}

The production win model is the map
\[
  \underbrace{(\vc_i,\ \Wp_i,\ \Ddiff_i)_{i=1}^N}_{\text{per-horse mechanism}}
   \;\longmapsto\;
  \Tfin_i\sim\IG\!\Big(\tfrac{\Leff_i}{\vc_i(1+\Wp_i/L)},\ \tfrac{L^2}{2\Ddiff_i}\Big)
   \;\longmapsto\;
  p_i = \Pr[\Tfin_i = \min_j \Tfin_j].
\]
The remainder of the document is devoted to the first arrow: how each of the three per-horse
parameters is built, leak-free, from data and mechanism. The second arrow is the Wald law of this
chapter; the third is the Monte-Carlo argmin of Chapter~\ref{ch:assembly}.
